
\documentclass{article}
\usepackage{colortbl}
\usepackage{makecell}
\usepackage{multirow}
\usepackage{supertabular}

\begin{document}

\newcounter{utterance}

\centering \large Interaction Transcript for game `dond', experiment `coop\_it', episode 18 with llama{-}3.3{-}70b{-}instruct{-}t1.0.
\vspace{24pt}

{ \footnotesize  \setcounter{utterance}{1}
\setlength{\tabcolsep}{0pt}
\begin{supertabular}{c@{$\;$}|p{.15\linewidth}@{}p{.15\linewidth}p{.15\linewidth}p{.15\linewidth}p{.15\linewidth}p{.15\linewidth}}
    \# & $\;$A & \multicolumn{4}{c}{Game Master} & $\;\:$B\\
    \hline

    \theutterance \stepcounter{utterance}  
    & & \multicolumn{4}{p{0.6\linewidth}}{
        \cellcolor[rgb]{0.9,0.9,0.9}{
            \makecell[{{p{\linewidth}}}]{
                \texttt{\tiny{[P1$\langle$GM]}}
                \texttt{State giocando a un gioco di negoziazione cooperativa in cui dovete accordarvi su come dividere una serie di oggetti tra voi e un altro giocatore.} \\
\\ 
\texttt{Regole:} \\
\texttt{(a) A Lei e all'altro giocatore viene dato un insieme di oggetti. Ognuno di voi riceve anche una funzione di valore segreta, che rappresenta il valore di ciascun tipo di oggetto.} \\
\texttt{(b) Si scambiano messaggi con l'altro giocatore per concordare chi si aggiudica gli oggetti. Potete inviare un massimo di 5 messaggi ciascuno, oppure terminare in anticipo facendo una proposta segreta in qualsiasi momento.} \\
\texttt{(c) A ciascuno di voi viene chiesto di inviare una proposta segreta indicando gli oggetti che desiderate, formattata tra parentesi quadre come segue: "{[}Proposta: <numero> <nome oggetto>, <numero> <nome oggetto>, <...>{]}".} \\
\texttt{(d) Se le vostre proposte sono complementari, cioè ci sono abbastanza oggetti per soddisfare entrambe le proposte, a ciascun giocatore viene assegnato un punteggio basato sulla somma dei valori degli oggetti ricevuti. In caso contrario, entrambi ricevono zero punti.} \\
\texttt{(e) L'obiettivo è massimizzare la somma del proprio punteggio e di quello dell'altro giocatore.} \\
\\ 
\texttt{Cominciamo.} \\
\\ 
\texttt{L'insieme degli oggetti disponibili è:} \\
\\ 
\texttt{1 moneta, 2 buste, 2 patate} \\
\\ 
\texttt{La funzione valore segreta è:} \\
\\ 
\texttt{moneta: 2, busta: 2, patata: 2} \\
\\ 
\texttt{IMPORTANTE: i vostri messaggi, a meno che non si tratti di una proposta segreta, vengono trasmessi direttamente all'altro giocatore, quindi non includete alcuna risposta alle regole o testo di annuncio del vostro messaggio. Per fare una proposta segreta, utilizzate il formato indicato. Non utilizzare delle parentesi quadre quando si comunica all'altro giocatore, altrimenti verrà interpretata come una proposta segreta.} \\
            }
        }
    }
    & \\ \\

    \theutterance \stepcounter{utterance}  
    & \multicolumn{4}{p{0.6\linewidth}}{
        \cellcolor[rgb]{0.9,0.9,0.9}{
            \makecell[{{p{\linewidth}}}]{
                \texttt{\tiny{[P1$\rangle$GM]}}
                \texttt{Ciao, vorrei iniziare proponendo di dividere gli oggetti in modo equo, considerando il valore e la quantità. Che ne pensi di prendere 1 busta e 1 patata, mentre io prendo l'altra busta, l'altra patata e la moneta?} \\
            }
        }
    }
    & & \\ \\

    \theutterance \stepcounter{utterance}  
    & & & \multicolumn{2}{p{0.3\linewidth}}{
        \cellcolor[rgb]{0.9,0.9,0.9}{
            \makecell[{{p{\linewidth}}}]{
                \texttt{\tiny{[GM$|$GM]}}
                \texttt{continue} \\
            }
        }
    }
    & & \\ \\

    \theutterance \stepcounter{utterance}  
    & & & \multicolumn{2}{p{0.3\linewidth}}{
        \cellcolor[rgb]{0.9,0.9,0.9}{
            \makecell[{{p{\linewidth}}}]{
                \texttt{\tiny{[GM$|$GM]}}
                \texttt{Ciao, vorrei iniziare proponendo di dividere gli oggetti in modo equo, considerando il valore e la quantità. Che ne pensi di prendere 1 busta e 1 patata, mentre io prendo l'altra busta, l'altra patata e la moneta?} \\
            }
        }
    }
    & & \\ \\

    \theutterance \stepcounter{utterance}  
    & & \multicolumn{4}{p{0.6\linewidth}}{
        \cellcolor[rgb]{0.9,0.9,0.9}{
            \makecell[{{p{\linewidth}}}]{
                \texttt{\tiny{[GM$\rangle$P2]}}
                \texttt{State giocando a un gioco di negoziazione cooperativa in cui dovete accordarvi su come dividere una serie di oggetti tra voi e un altro giocatore.} \\
\\ 
\texttt{Regole:} \\
\texttt{(a) A Lei e all'altro giocatore viene dato un insieme di oggetti. Ognuno di voi riceve anche una funzione di valore segreta, che rappresenta il valore di ciascun tipo di oggetto.} \\
\texttt{(b) Si scambiano messaggi con l'altro giocatore per concordare chi si aggiudica gli oggetti. Potete inviare un massimo di 5 messaggi ciascuno, oppure terminare in anticipo facendo una proposta segreta in qualsiasi momento.} \\
\texttt{(c) A ciascuno di voi viene chiesto di inviare una proposta segreta indicando gli oggetti che desiderate, formattata tra parentesi quadre come segue: "{[}Proposta: <numero> <nome oggetto>, <numero> <nome oggetto>, <...>{]}".} \\
\texttt{(d) Se le vostre proposte sono complementari, cioè ci sono abbastanza oggetti per soddisfare entrambe le proposte, a ciascun giocatore viene assegnato un punteggio basato sulla somma dei valori degli oggetti ricevuti. In caso contrario, entrambi ricevono zero punti.} \\
\texttt{(e) L'obiettivo è massimizzare la somma del proprio punteggio e di quello dell'altro giocatore.} \\
\\ 
\texttt{Cominciamo.} \\
\\ 
\texttt{L'insieme degli oggetti disponibili è:} \\
\\ 
\texttt{1 moneta, 2 buste, 2 patate} \\
\\ 
\texttt{La funzione valore segreta è:} \\
\\ 
\texttt{moneta: 2, busta: 1, patata: 3} \\
\\ 
\texttt{IMPORTANTE: i vostri messaggi, a meno che non si tratti di una proposta segreta, vengono trasmessi direttamente all'altro giocatore, quindi non includete alcuna risposta alle regole o testo di annuncio del vostro messaggio. Per fare una proposta segreta, utilizzate il formato indicato. Non utilizzare delle parentesi quadre quando si comunica all'altro giocatore, altrimenti verrà interpretata come una proposta segreta.} \\
\\ 
\\ 
\texttt{Ciao, vorrei iniziare proponendo di dividere gli oggetti in modo equo, considerando il valore e la quantità. Che ne pensi di prendere 1 busta e 1 patata, mentre io prendo l'altra busta, l'altra patata e la moneta?} \\
            }
        }
    }
    & \\ \\

    \theutterance \stepcounter{utterance}  
    & & & \multicolumn{4}{p{0.6\linewidth}}{
        \cellcolor[rgb]{0.9,0.9,0.9}{
            \makecell[{{p{\linewidth}}}]{
                \texttt{\tiny{[GM$\langle$P2]}}
                \texttt{Mi sembra una proposta interessante, ma prima di prendere una decisione vorrei discutere un po' di più sul valore degli oggetti. Non è necessariamente vero che la divisione che hai proposto sia la più equa, considerando che non conosciamo a vicenda i nostri valori segreti. Non potremmo trovare una divisione che massimizzi il valore totale per entrambi? Ad esempio, se uno di noi attribuisce un valore molto alto a un particolare oggetto, potrebbe essere meglio assegnarglielo e trovare un equilibrio con gli altri oggetti. Che idea ti sei fatto del valore degli oggetti?} \\
            }
        }
    }
     \\ \\

    \theutterance \stepcounter{utterance}  
    & & & \multicolumn{2}{p{0.3\linewidth}}{
        \cellcolor[rgb]{0.9,0.9,0.9}{
            \makecell[{{p{\linewidth}}}]{
                \texttt{\tiny{[GM$|$GM]}}
                \texttt{continue} \\
            }
        }
    }
    & & \\ \\

    \theutterance \stepcounter{utterance}  
    & & & \multicolumn{2}{p{0.3\linewidth}}{
        \cellcolor[rgb]{0.9,0.9,0.9}{
            \makecell[{{p{\linewidth}}}]{
                \texttt{\tiny{[GM$|$GM]}}
                \texttt{Mi sembra una proposta interessante, ma prima di prendere una decisione vorrei discutere un po' di più sul valore degli oggetti. Non è necessariamente vero che la divisione che hai proposto sia la più equa, considerando che non conosciamo a vicenda i nostri valori segreti. Non potremmo trovare una divisione che massimizzi il valore totale per entrambi? Ad esempio, se uno di noi attribuisce un valore molto alto a un particolare oggetto, potrebbe essere meglio assegnarglielo e trovare un equilibrio con gli altri oggetti. Che idea ti sei fatto del valore degli oggetti?} \\
            }
        }
    }
    & & \\ \\

    \theutterance \stepcounter{utterance}  
    & & \multicolumn{4}{p{0.6\linewidth}}{
        \cellcolor[rgb]{0.9,0.9,0.9}{
            \makecell[{{p{\linewidth}}}]{
                \texttt{\tiny{[P1$\langle$GM]}}
                \texttt{Mi sembra una proposta interessante, ma prima di prendere una decisione vorrei discutere un po' di più sul valore degli oggetti. Non è necessariamente vero che la divisione che hai proposto sia la più equa, considerando che non conosciamo a vicenda i nostri valori segreti. Non potremmo trovare una divisione che massimizzi il valore totale per entrambi? Ad esempio, se uno di noi attribuisce un valore molto alto a un particolare oggetto, potrebbe essere meglio assegnarglielo e trovare un equilibrio con gli altri oggetti. Che idea ti sei fatto del valore degli oggetti?} \\
            }
        }
    }
    & \\ \\

    \theutterance \stepcounter{utterance}  
    & \multicolumn{4}{p{0.6\linewidth}}{
        \cellcolor[rgb]{0.9,0.9,0.9}{
            \makecell[{{p{\linewidth}}}]{
                \texttt{\tiny{[P1$\rangle$GM]}}
                \texttt{Hai ragione, la divisione che ho proposto inizialmente potrebbe non essere ottimale. Considerando che i valori segreti possono variare, sarebbe utile scoprire se ci sono oggetti che hanno un valore particolarmente alto per uno di noi. Io personalmente attribuisco un valore piuttosto alto alla moneta, ma non so se questo sia vero anche per te. Forse potremmo cercare di scambiare informazioni senza rivelare troppo sui nostri valori segreti? Ad esempio, potremmo discutere se ci sono oggetti che uno di noi non è disposto a rinunciare o se ci sono oggetti che consideriamo più importanti degli altri senza necessariamente rivelare il nostro valore esatto.} \\
            }
        }
    }
    & & \\ \\

    \theutterance \stepcounter{utterance}  
    & & & \multicolumn{2}{p{0.3\linewidth}}{
        \cellcolor[rgb]{0.9,0.9,0.9}{
            \makecell[{{p{\linewidth}}}]{
                \texttt{\tiny{[GM$|$GM]}}
                \texttt{continue} \\
            }
        }
    }
    & & \\ \\

    \theutterance \stepcounter{utterance}  
    & & & \multicolumn{2}{p{0.3\linewidth}}{
        \cellcolor[rgb]{0.9,0.9,0.9}{
            \makecell[{{p{\linewidth}}}]{
                \texttt{\tiny{[GM$|$GM]}}
                \texttt{Hai ragione, la divisione che ho proposto inizialmente potrebbe non essere ottimale. Considerando che i valori segreti possono variare, sarebbe utile scoprire se ci sono oggetti che hanno un valore particolarmente alto per uno di noi. Io personalmente attribuisco un valore piuttosto alto alla moneta, ma non so se questo sia vero anche per te. Forse potremmo cercare di scambiare informazioni senza rivelare troppo sui nostri valori segreti? Ad esempio, potremmo discutere se ci sono oggetti che uno di noi non è disposto a rinunciare o se ci sono oggetti che consideriamo più importanti degli altri senza necessariamente rivelare il nostro valore esatto.} \\
            }
        }
    }
    & & \\ \\

    \theutterance \stepcounter{utterance}  
    & & \multicolumn{4}{p{0.6\linewidth}}{
        \cellcolor[rgb]{0.9,0.9,0.9}{
            \makecell[{{p{\linewidth}}}]{
                \texttt{\tiny{[GM$\rangle$P2]}}
                \texttt{Hai ragione, la divisione che ho proposto inizialmente potrebbe non essere ottimale. Considerando che i valori segreti possono variare, sarebbe utile scoprire se ci sono oggetti che hanno un valore particolarmente alto per uno di noi. Io personalmente attribuisco un valore piuttosto alto alla moneta, ma non so se questo sia vero anche per te. Forse potremmo cercare di scambiare informazioni senza rivelare troppo sui nostri valori segreti? Ad esempio, potremmo discutere se ci sono oggetti che uno di noi non è disposto a rinunciare o se ci sono oggetti che consideriamo più importanti degli altri senza necessariamente rivelare il nostro valore esatto.} \\
            }
        }
    }
    & \\ \\

    \theutterance \stepcounter{utterance}  
    & & & \multicolumn{4}{p{0.6\linewidth}}{
        \cellcolor[rgb]{0.9,0.9,0.9}{
            \makecell[{{p{\linewidth}}}]{
                \texttt{\tiny{[GM$\langle$P2]}}
                \texttt{Sì, è una buona strategia cercare di capire quali oggetti sono più importanti per l'altro senza rivelare troppo sui nostri valori segreti. Io posso dirti che le patate sembrano essere oggetti interessanti, ma non sono disposto a fare una mossa senza capire meglio la tua prospettiva. La moneta, come hai detto, potrebbe essere un punto di partenza interessante per la negoziazione. Sei disposto a considerare la possibilità di cedere la moneta in cambio di qualcos'altro, o è assolutamente fuori discussione?} \\
            }
        }
    }
     \\ \\

    \theutterance \stepcounter{utterance}  
    & & & \multicolumn{2}{p{0.3\linewidth}}{
        \cellcolor[rgb]{0.9,0.9,0.9}{
            \makecell[{{p{\linewidth}}}]{
                \texttt{\tiny{[GM$|$GM]}}
                \texttt{continue} \\
            }
        }
    }
    & & \\ \\

    \theutterance \stepcounter{utterance}  
    & & & \multicolumn{2}{p{0.3\linewidth}}{
        \cellcolor[rgb]{0.9,0.9,0.9}{
            \makecell[{{p{\linewidth}}}]{
                \texttt{\tiny{[GM$|$GM]}}
                \texttt{Sì, è una buona strategia cercare di capire quali oggetti sono più importanti per l'altro senza rivelare troppo sui nostri valori segreti. Io posso dirti che le patate sembrano essere oggetti interessanti, ma non sono disposto a fare una mossa senza capire meglio la tua prospettiva. La moneta, come hai detto, potrebbe essere un punto di partenza interessante per la negoziazione. Sei disposto a considerare la possibilità di cedere la moneta in cambio di qualcos'altro, o è assolutamente fuori discussione?} \\
            }
        }
    }
    & & \\ \\

    \theutterance \stepcounter{utterance}  
    & & \multicolumn{4}{p{0.6\linewidth}}{
        \cellcolor[rgb]{0.9,0.9,0.9}{
            \makecell[{{p{\linewidth}}}]{
                \texttt{\tiny{[P1$\langle$GM]}}
                \texttt{Sì, è una buona strategia cercare di capire quali oggetti sono più importanti per l'altro senza rivelare troppo sui nostri valori segreti. Io posso dirti che le patate sembrano essere oggetti interessanti, ma non sono disposto a fare una mossa senza capire meglio la tua prospettiva. La moneta, come hai detto, potrebbe essere un punto di partenza interessante per la negoziazione. Sei disposto a considerare la possibilità di cedere la moneta in cambio di qualcos'altro, o è assolutamente fuori discussione?} \\
            }
        }
    }
    & \\ \\

    \theutterance \stepcounter{utterance}  
    & \multicolumn{4}{p{0.6\linewidth}}{
        \cellcolor[rgb]{0.9,0.9,0.9}{
            \makecell[{{p{\linewidth}}}]{
                \texttt{\tiny{[P1$\rangle$GM]}}
                \texttt{Le patate sono interessanti, eh? Questo mi fa pensare che potresti attribuire loro un valore significativo. Quanto alla moneta, non sono assolutamente contrario a cederla, ma dipenderebbe da cosa potrei ottenere in cambio. Se potessi ottenere entrambe le buste, ad esempio, potrei considerare di cedere la moneta. Ma se sto a pensare a cosa potrei ottenere in cambio, mi viene in mente che forse potremmo fare una proposta che coinvolga gli oggetti in modo più creativo. Ad esempio, potresti prendere la moneta e una patata, mentre io prendo le due buste e l'altra patata? Sarebbe una divisione che tiene conto di più variabili e potrebbe massimizzare il valore per entrambi. Cosa ne pensi?} \\
            }
        }
    }
    & & \\ \\

    \theutterance \stepcounter{utterance}  
    & & & \multicolumn{2}{p{0.3\linewidth}}{
        \cellcolor[rgb]{0.9,0.9,0.9}{
            \makecell[{{p{\linewidth}}}]{
                \texttt{\tiny{[GM$|$GM]}}
                \texttt{continue} \\
            }
        }
    }
    & & \\ \\

    \theutterance \stepcounter{utterance}  
    & & & \multicolumn{2}{p{0.3\linewidth}}{
        \cellcolor[rgb]{0.9,0.9,0.9}{
            \makecell[{{p{\linewidth}}}]{
                \texttt{\tiny{[GM$|$GM]}}
                \texttt{Le patate sono interessanti, eh? Questo mi fa pensare che potresti attribuire loro un valore significativo. Quanto alla moneta, non sono assolutamente contrario a cederla, ma dipenderebbe da cosa potrei ottenere in cambio. Se potessi ottenere entrambe le buste, ad esempio, potrei considerare di cedere la moneta. Ma se sto a pensare a cosa potrei ottenere in cambio, mi viene in mente che forse potremmo fare una proposta che coinvolga gli oggetti in modo più creativo. Ad esempio, potresti prendere la moneta e una patata, mentre io prendo le due buste e l'altra patata? Sarebbe una divisione che tiene conto di più variabili e potrebbe massimizzare il valore per entrambi. Cosa ne pensi?} \\
            }
        }
    }
    & & \\ \\

    \theutterance \stepcounter{utterance}  
    & & \multicolumn{4}{p{0.6\linewidth}}{
        \cellcolor[rgb]{0.9,0.9,0.9}{
            \makecell[{{p{\linewidth}}}]{
                \texttt{\tiny{[GM$\rangle$P2]}}
                \texttt{Le patate sono interessanti, eh? Questo mi fa pensare che potresti attribuire loro un valore significativo. Quanto alla moneta, non sono assolutamente contrario a cederla, ma dipenderebbe da cosa potrei ottenere in cambio. Se potessi ottenere entrambe le buste, ad esempio, potrei considerare di cedere la moneta. Ma se sto a pensare a cosa potrei ottenere in cambio, mi viene in mente che forse potremmo fare una proposta che coinvolga gli oggetti in modo più creativo. Ad esempio, potresti prendere la moneta e una patata, mentre io prendo le due buste e l'altra patata? Sarebbe una divisione che tiene conto di più variabili e potrebbe massimizzare il valore per entrambi. Cosa ne pensi?} \\
            }
        }
    }
    & \\ \\

    \theutterance \stepcounter{utterance}  
    & & & \multicolumn{4}{p{0.6\linewidth}}{
        \cellcolor[rgb]{0.9,0.9,0.9}{
            \makecell[{{p{\linewidth}}}]{
                \texttt{\tiny{[GM$\langle$P2]}}
                \texttt{Questa proposta è più intrigante! Mi piace l'idea di una divisione che consideri più variabili e potenzialmente massimizzi il valore per entrambi. Tuttavia, prima di accettare, vorrei riflettere un po' sulla situazione. Se io prendo la moneta e una patata, mentre tu prendi le due buste e l'altra patata, potrei essere soddisfatto, ma dipende molto da come valuto realmente gli oggetti. Una cosa che potrei chiederti è: sei disposto a fare una piccola concessione o a provare a trovare un altro equilibrio, nel caso in cui la mia valutazione degli oggetti sia diversa da quella che tu ti aspetti? Ad esempio, potremmo scambiare una delle buste con qualcos'altro, o trovare un'altra combinazione che ci soddisfi entrambi? {[}Proposta: 1 moneta, 1 patata, 1 busta{]}} \\
            }
        }
    }
     \\ \\

    \theutterance \stepcounter{utterance}  
    & & & \multicolumn{2}{p{0.3\linewidth}}{
        \cellcolor[rgb]{0.9,0.9,0.9}{
            \makecell[{{p{\linewidth}}}]{
                \texttt{\tiny{[GM$|$GM]}}
                \texttt{proposal} \\
            }
        }
    }
    & & \\ \\

    \theutterance \stepcounter{utterance}  
    & & & \multicolumn{2}{p{0.3\linewidth}}{
        \cellcolor[rgb]{0.9,0.9,0.9}{
            \makecell[{{p{\linewidth}}}]{
                \texttt{\tiny{[GM$|$GM]}}
                [1, 1, 1]
            }
        }
    }
    & & \\ \\

    \theutterance \stepcounter{utterance}  
    & & \multicolumn{4}{p{0.6\linewidth}}{
        \cellcolor[rgb]{0.9,0.9,0.9}{
            \makecell[{{p{\linewidth}}}]{
                \texttt{\tiny{[P1$\langle$GM]}}
                \texttt{L'altro giocatore ha presentato una proposta segreta. Ora tocca a Lei presentare una sua proposta.} \\
            }
        }
    }
    & \\ \\

    \theutterance \stepcounter{utterance}  
    & \multicolumn{4}{p{0.6\linewidth}}{
        \cellcolor[rgb]{0.9,0.9,0.9}{
            \makecell[{{p{\linewidth}}}]{
                \texttt{\tiny{[P1$\rangle$GM]}}
                \texttt{{[}Proposta: 1 moneta, 1 busta, 1 patata{]}} \\
            }
        }
    }
    & & \\ \\

    \theutterance \stepcounter{utterance}  
    & & & \multicolumn{2}{p{0.3\linewidth}}{
        \cellcolor[rgb]{0.9,0.9,0.9}{
            \makecell[{{p{\linewidth}}}]{
                \texttt{\tiny{[GM$|$GM]}}
                \texttt{proposal} \\
            }
        }
    }
    & & \\ \\

    \theutterance \stepcounter{utterance}  
    & & & \multicolumn{2}{p{0.3\linewidth}}{
        \cellcolor[rgb]{0.9,0.9,0.9}{
            \makecell[{{p{\linewidth}}}]{
                \texttt{\tiny{[GM$|$GM]}}
                [1, 1, 1]
            }
        }
    }
    & & \\ \\

    \theutterance \stepcounter{utterance}  
    & & & \multicolumn{2}{p{0.3\linewidth}}{
        \cellcolor[rgb]{0.9,0.9,0.9}{
            \makecell[{{p{\linewidth}}}]{
                \texttt{\tiny{[GM$|$GM]}}
                [[1, 1, 1], [1, 1, 1]]
            }
        }
    }
    & & \\ \\

\end{supertabular}
}

\end{document}
